Tool-description quality is widely treated as a broadly-applicable lever
for agent tool-use, but it is not a single better/worse axis: the
precision that helps an agent disambiguate within a family of confusable
tools is orthogonal to, or actively harmful for, context-rich selection
and for tool retrieval. We test this with a single frozen evaluation
protocol --- one classifier, one judge, one generator family,
pre-registered thresholds --- across a synthetic confusable-catalog
experiment, two real production MCP-server mirrors (GitHub, AWS IAM), a
synthetic internal-proxy catalog, and a pre-registered pilot of ten
public Python MCP servers. The effect is real but regime-bounded, not a
general law: an oracle description helps at one tested catalog density
(60 tools/10 families, +34.5pp on gemma2:9b, and not collapsing on a
substantially stronger model, Llama-3.3-70B) with no headroom; realizing
it safely through automatic generation is a separate condition,
requiring documented source. Outside these conditions the effect is
null or reverses (zero headroom on two well-documented production
servers; −20pp harm on one already-resolved family; harm to retrieval
across three retriever types). Contributions: a falsifiable regime map
of where description quality helps, harms, or does nothing; a
pre-registered prevalence measurement finding the behavioral regime in
0 of 9 testable Python MCP servers --- a lower bound, not a population
estimate; and a localizability boundary --- a pairwise LLM-judge
confusability method fails under two independent framings, via two
distinct failure mechanisms.
